%%%===========================================
% \phantomsection
% \addcontentsline{toc}{section}{Abstract}

% =============================
% \section*{Abstract}
% =============================

Environmental degradation is often associated with a structural trade-off between increasing economic productivity and preserving the ecological systems that sustain production over time. Although economic theory has long considered nature through concepts such as land, natural resources, externalities, environmental costs, and common-pool resources, conventional production theory still tends to treat the environment primarily as a constraint, stock, or corrective problem rather than as an endogenous, dynamic, and strategically relevant component of production. This separation contributes to institutional arrangements in which preservation is frequently perceived by private agents as an external obligation, a compliance cost, or a restriction on productive activity.

This research proposes an alternative theoretical framework based on the concept of \emph{economic commensalism}. The concept is used to describe an institutional and economic arrangement in which environmental preservation is not treated merely as a moral duty, legal constraint, or compensatory mechanism, but as a condition that can be incorporated into the agent's productive calculus. The central argument is that preservation may become individually rational when ecological variables are endogenized into production decisions and when incentives are spatially, intertemporally, and strategically calibrated. In this sense, the research moves beyond approaches that rely exclusively on Pigouvian taxation, Coasian bargaining, command-and-control regulation, or fragmented positive incentives. Instead, it investigates whether environmental preservation can alter the structural payoff of economic agents and become compatible with private productivity, profitability, and long-term welfare.

The thesis is structured around three complementary components. First, it develops a theoretical framework that revisits the relationship between production, preservation, extractivism, predation, and common-pool resource governance. Second, it proposes a formal model in which environmental variables, intertemporal decision-making, uncertainty, bounded rationality, and strategic interaction are incorporated into the analysis of agents' choices. Third, it evaluates the empirical feasibility of the proposed framework by applying it to the Brazilian state of Mato Grosso, a territory marked by agricultural expansion, deforestation pressures, and ecological heterogeneity across biomes.

The research therefore seeks to contribute to environmental and resource economics by reframing environmental preservation as a potentially endogenous component of economic rationality. Its main contribution is to examine whether locally grounded and incentive-compatible governance mechanisms can reduce the opposition between production and preservation, making ecological conservation not only socially desirable, but also strategically stable and economically rational for private agents.

\newpage

% fim %
